\section{Introduction}
\label{sec:intro}

Accredited testing laboratories must review large volumes of conformity assessment documents under strict traceability requirements. In China National Accreditation Service for Conformity Assessment (CNAS) accredited settings, each test report is checked against accreditation criteria, test standards, and internal standard operating procedures, covering format validity, numerical consistency, conclusion logic, and normative references. At our partner laboratory, manual review takes about 35 minutes per report, varies across reviewers, and becomes increasingly costly as regulations change.

Large language models (LLMs) offer strong semantic understanding but remain risky in compliance-sensitive workflows because of hallucination \cite{ji2023hallucination,huang2023raghallucination}. Retrieval-augmented generation (RAG) grounds generation in external corpora \cite{lewis2020rag,gao2024ragsurvey}, yet retrieval failure, knowledge conflict, and reasoning drift can still yield unsupported claims \cite{gao2024ragsurvey,huang2023raghallucination}. Rule-based automated compliance checking remains auditable and effective on formally encoded requirements, but semantic regulatory text still calls for hybrid semantic--rule pipelines \cite{solihin2015rules,zhang2016acc,zheng2022acc}. Existing ACC and RAG systems seldom specify when priority routing should replace result fusion, or how reproducible mechanism tests should be separated from deployment-only evidence.

We propose \HRAG{}-DocAudit, a priority-cascaded hybrid RAG framework for document compliance auditing. When keyword rule signals fire, the system follows deterministic rule decisions; remaining items are handled by document-scoped hybrid retrieval (dense plus sparse) and LLM judgment, with confidence-tiered human-in-the-loop (HITL) escalation in deployment \cite{shneiderman2020hcai}. We make three contributions: (i)~a priority cascade that isolates RAG-origin unsupported positives on rule-routed items (lenient metric) while preserving strict false-positive parity with OR-ensemble, evaluated with stratified public-track tests; (ii)~document-scoped hybrid retrieval with a validation $\lambda$ ablation that separates ranking (NDCG@10) from classification accuracy; and (iii)~confidence-tiered HITL governance that lowers auditor load under a CNAS partner-laboratory deployment. A dual-track protocol supports these claims: Tier-1 open replication on C3PA and ContractNLI isolates routing and retrieval, while Tier-2 CNAS deployment reports end-to-end accuracy, lenient hallucination, HITL load, and latency under accredited-laboratory constraints.
